\lampiransection{HASIL PRA-SURVEI TAHAP 2}
\label{app:prasurvey2}

Lampiran ini memuat hasil pra-survei tahap 2 penelitian \textit{Analisis Strategi Pemasaran Rumah Makan Nasi Gerilya pada Platform GrabFood Menggunakan Matriks SWOT}. Pra-survei tahap 2 dilakukan melalui wawancara awal dengan pemilik usaha untuk memperoleh data mengenai profil usaha, penggunaan platform, struktur pelanggan, kekuatan usaha, strategi pemasaran, kendala operasional, peluang pengembangan, dan ancaman persaingan. Seluruh poin pada lampiran ini disusun berdasarkan data dan keterangan yang disampaikan oleh pemilik usaha pada saat wawancara awal.

\section{Profil Usaha, Platform, dan Pelanggan}

\begin{table}[h!]
    \centering
    \caption{Profil Umum Usaha Berdasarkan Wawancara Pemilik}
    \label{tab:profil_prasurvey2}
    \begin{tabularx}{\textwidth}{|l|X|}
        \hline
        \textbf{Aspek}             & \textbf{Keterangan}              \\
        \hline
        Nama usaha                 & Rumah Makan Nasi Gerilya         \\
        \hline
        Tahun berdiri              & September 2024                   \\
        \hline
        Lama usaha saat pra-survei & Sekitar 1 tahun 7 bulan          \\
        \hline
        Kanal penjualan online     & GrabFood, GoFood, dan ShopeeFood \\
        \hline
    \end{tabularx}
\end{table}

\begin{table}[h!]
    \centering
    \caption{Keterangan Pemilik Mengenai Penggunaan Platform}
    \label{tab:prasurvey2_grabfood}
    \begin{tabularx}{\textwidth}{|l|X|}
        \hline
        \textbf{Aspek}                & \textbf{Keterangan Berdasarkan Wawancara}                                                                       \\
        \hline
        Platform utama                & GrabFood                                                                                                        \\
        \hline
        Kontribusi terhadap penjualan & Disebut sebagai penyumbang penjualan terbesar                                                                   \\
        \hline
        Pertimbangan pemilik          & Segmen pengguna GrabFood dinilai lebih sesuai dengan target pasar usaha                                         \\
        \hline
        Dukungan platform             & Awards, \textit{verified store}, dan \textit{voucher discount}                                                  \\
        \hline
        Pandangan terhadap GoFood     & Dinilai lebih banyak memiliki sebaran pengguna di pinggiran kota dan kurang sesuai dengan segmen pasar usaha    \\
        \hline
        Pandangan terhadap ShopeeFood & Dinilai lebih banyak mendorong pembelian berbasis voucher; pembelian tidak selalu berujung pada pembelian ulang \\
        \hline
    \end{tabularx}
\end{table}

\begin{table}[h!]
    \centering
    \caption{Komposisi Pelanggan Berdasarkan Keterangan Pemilik}
    \label{tab:prasurvey2_pelanggan}
    \begin{tabularx}{\textwidth}{|X|c|c|c|}
        \hline
        \textbf{Kategori Pelanggan} & \textbf{Periode Sebelumnya} & \textbf{Saat Ini} & \textbf{Perubahan} \\
        \hline
        \textit{New-comers}         & 45\%                        & 55\%              & +10 pp             \\
        \hline
        \textit{Repeat order}       & 43\%                        & 38\%              & -5 pp              \\
        \hline
        \textit{Reactivated user}   & 12\%                        & 7\%               & -5 pp              \\
        \hline
    \end{tabularx}
\end{table}

Tabel~\ref{tab:profil_prasurvey2}, Tabel~\ref{tab:prasurvey2_grabfood}, dan Tabel~\ref{tab:prasurvey2_pelanggan} menunjukkan bahwa pra-survei tahap 2 memberikan data awal mengenai identitas usaha, alasan pemilihan platform utama, serta perubahan komposisi pelanggan dari periode sebelumnya ke periode saat ini. Pemilik usaha mengaitkan perubahan komposisi pelanggan tersebut dengan pelaksanaan \textit{marketing plan high season}.

\begin{table}[h!]
    \centering
    \caption{Keterangan Pemilik Mengenai Kekuatan dan Strategi Usaha}
    \label{tab:kekuatan_strategi_prasurvey2}
    \begin{tabularx}{\textwidth}{|l|X|}
        \hline
        \textbf{Aspek}            & \textbf{Keterangan Berdasarkan Wawancara}                                                                                                      \\
        \hline
        Kekuatan utama            & Rasa produk dinilai otentik dan \textit{brand recognition} dinilai sudah cukup kuat                                                            \\
        \hline
        Produk unggulan           & Dendeng dan ayam pop putih                                                                                                                     \\
        \hline
        Strategi                  & \textit{Upselling} melalui optimisasi \textit{product lineup}; peningkatan kualitas produk; peningkatan layanan; \textit{discount partnership} \\
        \hline
        Peran GrabFood            & Digunakan lebih banyak sebagai kanal \textit{call to action}                                                                                   \\
        \hline
        Kanal pemasaran utama     & Media sosial serta KOL atau \textit{influencer}                                                                                                \\
        \hline
        Penguatan identitas merek & Perancangan maskot, \textit{signature gesture}, dan budaya merek yang unik                                                                     \\
        \hline
    \end{tabularx}
\end{table}

\begin{table}[h!]
    \centering
    \caption{Kendala Operasional Berdasarkan Wawancara Pemilik}
    \label{tab:kendala_prasurvey2}
    \begin{tabularx}{\textwidth}{|l|X|}
        \hline
        \textbf{Kategori Kendala}       & \textbf{Keterangan Berdasarkan Wawancara}                                                                                                                                                                                                                         \\
        \hline
        Kendala umum                    & Ongkos kirim dinilai mahal; jarak layanan masih terbatas; harga mengalami kenaikan karena usaha sedang meningkatkan margin dan \textit{net profit}; diskon, pajak, dan biaya promo tidak lagi sepenuhnya ditanggung platform                                      \\
        \hline
        Kendala layanan                 & Antrean driver sering terjadi terutama pada \textit{peak season}; pemilik menginginkan sistem antrean yang memprioritaskan driver yang datang lebih dahulu; makanan yang habis masih dapat dipesan; pesanan kadang tidak sesuai dengan \textit{request} pelanggan \\
        \hline
        \textit{Bottleneck} operasional & Pada \textit{peak hour} banyak menu \textit{sold out} dan perlu \textit{restock}; pelanggan dapat membatalkan pembelian ketika menu habis; kapasitas penggorengan masih kurang; belum tersedia opsi \textit{dessert} untuk mendukung \textit{upselling}           \\
        \hline
        Konsistensi rasa                & Pemilik menegaskan konsistensi rasa masih menjadi masalah utama; sebagian \textit{cook} masih mengambil \textit{shortcut}; diperlukan pengawasan yang lebih ketat agar rasa tetap konsisten                                                                       \\
        \hline
    \end{tabularx}
\end{table}

\begin{table}[h!]
    \centering
    \caption{Peluang dan Ancaman Berdasarkan Wawancara Pemilik}
    \label{tab:peluang_ancaman_prasurvey2}
    \begin{tabularx}{\textwidth}{|l|X|}
        \hline
        \textbf{Kategori}  & \textbf{Keterangan Berdasarkan Wawancara}                                                                                                                                                                                           \\
        \hline
        Peluang platform   & Dukungan program GrabFood seperti awards, \textit{verified store}, dan \textit{voucher discount}; peluang \textit{upselling} melalui pengembangan produk; potensi peningkatan nilai transaksi melalui \textit{discount partnership} \\
        \hline
        Peluang ekspansi   & Usaha belum memiliki cabang; pemilik ingin membuka cabang kedua pada tahun ini; terdapat peluang membuka cabang baru di Kota Medan; target jangka panjang adalah ekspansi ke Jakarta setelah memiliki lima cabang di Medan          \\
        \hline
        Ancaman persaingan & Munculnya kompetitor baru dengan \textit{brand recognition} kuat di Kota Medan, seperti Pagi Sore dan Payakumbuah                                                                                                                   \\
        \hline
        Ancaman platform   & Tingginya ongkos kirim; pergeseran beban promo kepada konsumen; sensitivitas pelanggan terhadap harga                                                                                                                               \\
        \hline
    \end{tabularx}
\end{table}

\begin{table}[h!]
    \centering
    \caption{Kategorisasi SWOT Awal Berdasarkan Wawancara Pemilik}
    \label{tab:swot_prasurvey2}
    \begin{tabularx}{\textwidth}{|l|X|}
        \hline
        \textbf{Kategori}      & \textbf{Data Awal}                                                                                                                                                                                                                                                                                           \\
        \hline
        \textit{Strengths}     & Rasa produk dinilai otentik; \textit{brand recognition} dinilai cukup kuat; GrabFood disebut sebagai penyumbang terbesar penjualan; terdapat produk unggulan berupa dendeng dan ayam pop putih; usaha mendapat dukungan program platform                                                                     \\
        \hline
        \textit{Weaknesses}    & Proporsi \textit{repeat order} menurun; konsistensi rasa belum stabil; makanan yang habis masih dapat dipesan; pesanan kadang tidak sesuai; terjadi \textit{bottleneck} operasional saat \textit{peak hour}; kapasitas penggorengan masih terbatas; belum tersedia \textit{dessert} untuk \textit{upselling} \\
        \hline
        \textit{Opportunities} & Terdapat peluang \textit{upselling}; peluang \textit{discount partnership}; dukungan fitur dan program promosi GrabFood; rencana penguatan identitas merek; peluang ekspansi cabang di Medan dan Jakarta                                                                                                     \\
        \hline
        \textit{Threats}       & Munculnya kompetitor baru dengan \textit{brand recognition} kuat; ongkos kirim dinilai mahal; jarak pengantaran terbatas; kenaikan harga berpotensi memengaruhi pembelian; beban promo semakin banyak ditanggung pengguna                                                                                    \\
        \hline
    \end{tabularx}
\end{table}

\section{Ringkasan dan Fungsi Pra-Survei Tahap 2}

Berdasarkan wawancara awal, pemilik usaha menyampaikan data mengenai posisi GrabFood sebagai platform utama, komposisi pelanggan, kekuatan produk, strategi pemasaran, kendala operasional, peluang pengembangan, dan ancaman eksternal. Seluruh data tersebut digunakan sebagai bahan awal untuk mengidentifikasi faktor internal dan eksternal usaha pada analisis SWOT penelitian.

Hasil pra-survei tahap 2 digunakan sebagai dasar untuk:
\begin{enumerate}[nosep]
    \item memperdalam identifikasi faktor internal dan eksternal usaha;
    \item melengkapi hasil observasi pada pra-survei tahap 1;
    \item menyusun pedoman wawancara penelitian utama;
    \item memperkuat dasar penyusunan matriks SWOT penelitian.
\end{enumerate}
